\section{Sharp bounds and the limiting tree measures}
\label{sec:sharp-asymptotics}

This section proves the estimate at the fluctuation scale.  The estimates
imported from \cite{DHM2025} are used through
\Cref{thm:marked-expansion} only to remove large complexity and expansion
remainders.  The sharper estimate for the diagrams that remain is proved
below by a change of relative translations. Sharp connected cumulants
and minimal-tree limits are central to the short-time fluctuation
theory of~\cite{BGSS2023}. The present section must establish the
uniform bounds and summability needed on the full regular kinetic
interval.

\subsection{Conventions and the assertion}

Fix a finite list of observation times and bounded continuous observation
functions.  When observations are accumulated in a mark, its update is
\(y\mapsto y+\phi_j(x,v)\), so that \(|y|\le M\), with \(M\) independent
of the number of particles.  More generally, one may retain a finite list
of phase points along each designated trajectory.  A record is retained
only along a line on which it is required.  Forgetting part of a record
means pushing the corresponding measure forward under its coordinate
projection.  This operation commutes with the finite hierarchy before
the independent factors are regrouped, by
\Cref{thm:marked-expansion}.

Let \(c_h^\eps\) denote the ordered, normalized factorial cumulant measure
of the pasted process, including these records on its \(h\) designated
particles.  Thus \(c_1^\eps\) is its one-particle measure, and the
normalization of an \(h\)-particle factorial measure is \(\mu^{-h}\).
Constants with a subscript \(K\) below may depend on a fixed upper bound
\(K\) for complexity, on the fixed observations, and on \(L\).  The
constant in an exponential factor \(C^n\) will not depend on \(K\).

The independent factors supplied by \Cref{prop:one-root} satisfy, for
some \(b>0\) and \(B_0<\infty\),
\begin{equation}\label{eq:sharp-input-bounds}
 \sup_{0<\eps\le\eps_0,\,i}\|g_i\|_{\Bol,b}\le B_0,
 \qquad
 \|g_i-f_{i\tau}\|_{\Bol,b}\longrightarrow0.
\end{equation}
We decrease \(b\) a fixed number of times when reserving Gaussian
weights.  This does not change the choice of the microscopic dynamics.

\begin{theorem}[Sharp connected bounds]\label{thm:sharp-cumulants}
Choose the parameters \(\Gamma,L\) as in
\Cref{thm:marked-expansion}, increasing \(L\), if necessary,
by an amount depending only on the uniform bounds for the Boltzmann
solution and on \(T\).  This one choice works for every fixed factorial
order and every fixed bounded observation list.  For each fixed \(h\),
\begin{equation}\label{eq:sharp-full-cumulant}
 \|c_h^\eps\|_{\TV}\le C_h\mu^{1-h}.
\end{equation}
For \(h=1,2\), the measures \(\mu^{h-1}c_h^\eps\) converge against
bounded continuous record tests to the signed sums of the retained
global trees with trivial initial cumulants.  At a contact their weight
is the incoming flux
\([(v-v_*)\cdot\omega]_+\,\dd t\,\dd\omega\); a scattering contact
has sign \(+1\), and a non-scattering contact has sign \(-1\).
The independent birth factors are \(f_{i\tau}\).

The sums of the absolute retained diagram measures, after multiplication
by \(\mu^{h-1}\), are bounded uniformly with a sufficiently small
Gaussian weight in their designated endpoint velocities.  The same
bound holds for the limiting tree measures, and for the Gaussian weights
in finitely many retained record velocities after decreasing the exponent.
Among the retained connected diagrams, the sum of the absolute measures
having a nontrivial initial cumulant or a global cycle is
\(o(\mu^{1-h})\).  The assertions about Gaussian weights concern retained
diagrams and their limits; no weighted estimate for a microscopic
expansion remainder is asserted.
\end{theorem}

Here and below a global tree is a tree on the particle lines, with a
contact as an edge.  Continuing one physical line through an interface
does not create a new vertex.  Different particles newly integrated at
an interface have different vertices.  This convention is essential for
counting both translations and activity factors.

\subsection{The finite graph and its numerical data}

Work first with a term of complexity at most \(K\) in the finite
expansion.  Replace each initial cumulant by its bound from
\Cref{prop:initial}, in the form
\begin{equation}\label{eq:sharp-initial-link}
 C_K\sum_{\mathcal L\text{ a forest}}
 \prod_{(p,q)\in\mathcal L}
 \left(\ind_{|x_p(0)-x_q(0)|\le R_{K,D}\eps}+\eps^D\right)
 \prod_p f_0(z_p(0)).
\end{equation}
Only the bounded set of vertices belonging to a nontrivial initial
cumulant occurs in \(\mathcal L\).  This replacement is used for absolute
bounds, not to assign a sign to an initial link.  Expand the parentheses
and, initially, suppose all chosen links are physical, namely the
indicator terms.

Let \(m\) be the number of distinct particle lines, \(c\) the number of
global components, \(p\) the number of physical initial links, and
\(e_i\) the number of contacts in slab \(i\).  The initial link graph is
a forest.  Start with its components and then examine contacts in time
order.  A contact is a \emph{merger} if its two lines belong to different
components immediately before that contact; otherwise call it redundant.
Write \(a_i\) for the number of mergers in slab \(i\), and \(r\) for
the number of redundant contacts.  The elementary forest identities are
\begin{equation}\label{eq:sharp-merger-count}
 \sum_i a_i=m-c-p,
 \qquad
 r=\sum_i e_i+p-m+c.
\end{equation}
The second expression is the global cycle excess.

Let \(n_i\) be the number of particle lines present in slab \(i\),
including lines with no contact there, and let \(q\) be the sum of the
numbers of lines crossing slab interfaces.  Then
\begin{equation}\label{eq:sharp-line-count}
 \sum_i n_i=m+q,
 \qquad
 \sum_i(n_i-a_i)=q+c+p.
\end{equation}
The merger edges in any single slab form a forest on its \(n_i\)
lines, so \(a_i\le n_i\).  All of \(q,p,c,r\), and hence
\(n_i-a_i\), are bounded in terms of \(K\).  To see the assertion for
\(r\) directly from the layer complexity, glue the layer graphs along
their interface lines: the total cycle excess can increase by at most
the number of gluings.  The layer excesses and these gluings are included
in \(\rho\).  Each layer component meets a layer root, so the number of
its components is bounded as well.  Erasing a bounded number of initial
links changes these assertions only by a bounded amount.

There is one independent input velocity and one input position for each
of the \(m\) birth lines.  A continued line has neither a new position
variable nor a new velocity variable.  The cancellation of activity
factors at interfaces, in the normalized factorial expansion, leaves
the factor
\begin{equation}\label{eq:sharp-activity}
                         \mu^{m-h}.
\end{equation}
In particular, the birth at an interface of a line supplied by \(g_i\)
is an independent phase variable, whereas the passage through that
interface of a recorded line is a continuation of its old variable.

\subsection{Changing the relative translations}

\begin{figure}[tb]
\centering
\begin{tikzpicture}[x=1cm,y=1cm,
  particle/.style={circle,draw=black,fill=white,minimum size=6mm,inner sep=1pt},
  merger/.style={thick},
  extra/.style={thick,dashed,draw=linkblue},
  lab/.style={font=\small,fill=white,inner sep=2pt}]
 \node[particle] (a) at (0,0) {$a$};
 \node[particle] (b) at (1.6,1) {$b$};
 \node[particle] (c) at (3.2,1) {$c$};
 \node[particle] (d) at (4.8,0) {$d$};
 \draw[merger] (a)--node[lab,above left] {$t_1$}(b);
 \draw[merger] (b)--node[lab,above] {$t_3$}(c);
 \draw[merger] (c)--node[lab,above right] {$t_2$}(d);
 \draw[extra] (a) to[bend right=24]
    node[lab,below] {$t_4$: redundant}(d);
 \node[anchor=west,font=\small,align=left] at (6.1,.95)
   {$t_1:\ \{a\}\cup\{b\}\longrightarrow\{a,b\}$\\[5pt]
    $t_2:\ \{c\}\cup\{d\}\longrightarrow\{c,d\}$\\[5pt]
    $t_3:\ \{a,b\}\cup\{c,d\}\longrightarrow\{a,b,c,d\}$};
 \node[font=\small] at (2.4,-1.25) {$t_1<t_2<t_3<t_4$};
\end{tikzpicture}
\caption{Chronological connectivity for four particle lines. Each solid
contact joins two previously separate components and fixes one relative
translation, with the incoming-flux Jacobian \(\eps^2B\,\dd t\,\dd\omega\).
The dashed contact is internal to the already joined component. Its
relative position is fixed, so it contributes a constraint rather than a
new translation variable. Later common translations cannot change that
constraint. Edges depict the contact graph, not physical trajectories.}
\label{fig:chronological-mergers}
\end{figure}

The distinction between a new merger and a redundant contact is
illustrated in \Cref{fig:chronological-mergers}.  It refers to the
chronological connected components of the contact graph, not to spatial
proximity at an unspecified time.

\begin{lemma}[Chronological merger coordinates]\label{lem:sharp-merger-coordinates}
Fix a finite line graph, the C/O choices, the chronological order, and
which contacts are redundant in that order.  Away from nontransversal
contacts and simultaneous events, its history integral has coordinates
consisting of the birth velocities, one absolute translation for each
global component, a vector in a ball of radius \(R_{K,D}\eps\) for each
physical initial merger, and
\((t,\omega)\in\R\times S^2\) for each dynamic merger.  The positional
Jacobian at a dynamic merger is
\begin{equation}\label{eq:sharp-jacobian}
       \eps^2[(v-v_*)\cdot\omega]_+\,\dd t\,\dd\omega.
\end{equation}
A specified redundant contact has no free continuous parameter and
requires no extra Jacobian.  The coordinate map has multiplicity at most
one on each prescribed incoming straight-segment branch.
\end{lemma}

\begin{proof}
Choose the birth velocities first.  Initially every line has an
independent translation of its free affine trajectory.  An initial
forest link determines the difference between the translations of its
two groups by a vector in the stated ball.  Orient the forest toward
one representative per group.  The changes from individual translations
to representatives and these differences have determinant one.

Proceed in increasing contact time.  Inductively, the histories within
each group have been constructed, including all velocity changes before
the next merger, and are known up to a common translation of that group.
For a merger between groups with translations \(X\) and \(X_*\), the
contact condition on the indicated straight segments is
\[
 X-X_*=-d(t)-\eps\omega,
\]
where \(d(t)\) is the known relative displacement with the translations
removed, and \(\omega\) is opposite to the outward geometric relative
normal.  Thus incoming contact has \((v-v_*)\cdot\omega>0\), as in
\eqref{eq:collision-kernel}.  On these segments
\(\dot d(t)=v-v_*\).  The determinant of
the map from time and surface coordinates on \(S^2\) to \(X-X_*\) is
\(\eps^2|(v-v_*)\cdot\omega|\).  The incoming convention selects the
sign in \eqref{eq:sharp-jacobian}.  A scattering contact applies the
elastic velocity reflection; an O contact retains the velocities.  The
two groups now have one common free translation.

If the next specified contact is redundant, its relative displacement
and velocities on the indicated segments are already fixed.  A straight
relative trajectory has at most one incoming intersection with the
sphere of radius \(\eps\).  More precisely, for its relative flight
\(r(t)=d+tu\) on the common segment interval \(I\),
\[
 \int_I\delta(|r(t)|-\eps)
       [-u\cdot r(t)/\eps]_+\,\dd t
 =\ind_{\{\text{an incoming intersection occurs in }I\}}.
\]
At a transversal incoming crossing, the flux cancels the absolute
derivative of \(|r(t)|-\eps\).  The outgoing delta factors then make
the prescribed C or O substitution.  Thus its time and normal, when
the prescribed intersection exists, are determined.  If it does not exist, the branch
has weight zero.  This includes segments starting inside the sphere:
their outgoing exit is not an incoming intersection.  Subsequent
translations of the entire group leave this test and its resulting
velocity reflection unchanged.

Every step therefore uses only coordinates fixed at earlier steps.  Its
derivatives with respect to those earlier coordinates may be nonzero,
but the full determinant is triangular and is the product of the
displayed diagonal determinants.  There is no new unknown relative
translation at a redundant step.  Contacts at different birth interfaces
are treated identically: the affine origins of the new lines differ,
but the derivative of their current displacement is still their current
velocity.  At the end retain one representative translation per global
component.  On an incoming transversal branch the reverse construction
recovers every relative translation, so the multiplicity is one.

For a finite graph the excluded segment boundaries, equality of two
specified times, and zero incoming flux have measure zero in merger
coordinates unless they belong to a redundant-contact requirement.
Such a requirement is retained as an indicator; no change of variables
is performed at it.  Equivalently, one may first work on compact subsets
of the transversal branches and then exhaust them.  This also justifies
the formula for nonnegative history integrals without smooth densities.
\end{proof}

Multiplying \eqref{eq:sharp-activity} by the merger Jacobians and the
physical link volumes gives
\begin{equation}\label{eq:sharp-scaling}
 \mu^{m-h}\eps^{2(m-c-p)}\eps^{3p}
                     =\mu^{c-h}\eps^p.
\end{equation}
All history acceptance indicators can be retained during the coordinate
change.  For an upper bound they can subsequently be dropped, except
that a redundant scattering is used to determine the velocities of its
specified branch when it occurs.  Dropping the requirement of its
occurrence is understood as extending the resulting integrand by zero
off that branch, not as assigning an arbitrary additional velocity.

For a connected diagram \(c=1\).  A nontrivial physical initial block
has \(p\ge1\), and thus gains at least one factor \(\eps\).  If a term
in \eqref{eq:sharp-initial-link} instead uses a tail factor, erase that
link.  Erasing at most \(K\) links increases the number of components
by at most \(K\), and the extra activity cost in
\eqref{eq:sharp-scaling} is at most \(\eps^{-2K}\).  Choosing
\(D>2K+M\) makes the terms with at least one such tail
\(O(\eps^M\mu^{1-h})\).  Here \(M\) is arbitrary and is chosen after
the required moment order.  The constants in a physical link radius
and the forest count involve only boundedly many links.

\subsection{Factorials and a summable majorant}

\begin{lemma}[The fresh-label factorials]\label{lem:sharp-factorials}
For diagrams of complexity at most \(K\), the absolute sum with slab
sizes \((n_1,\ldots,n_L)\) is bounded by a sum over chronologically
ordered merger pair lists with denominator \(\prod_i n_i!\), multiplied
by
\begin{equation}\label{eq:sharp-polynomial-count}
 C_K(1+n)^{d_K}C^{\sum_i a_i},
 \qquad n=\sum_i n_i,
\end{equation}
where the base \(C\) does not depend on \(K\).  The merger pairs
themselves are not counted in \eqref{eq:sharp-polynomial-count}.
\end{lemma}

\begin{proof}
In slab \(i\), let \(s_i\) be the number of labels arriving as roots
from above, including the final designated roots in the highest slab.
The finite factorial expansion integrates \(j_i=n_i-s_i\) new labels
with coefficient \(1/j_i!\).  Since \(s_i\le K\),
\[
 \frac1{j_i!}=\frac{n_i!}{(n_i-s_i)!}\frac1{n_i!}
             \le\frac{(1+n_i)^{s_i}}{n_i!}.
\]
The choices of continued labels, dependent subsets, and labels in
initial links all involve boundedly many labels; their number is a
fixed power of \(1+n\).

The rooted compatibility expansion in \Cref{thm:marked-expansion}
represents each connected cluster coefficient by spanning forests.
Its deterministic exploration gives at most one branch per selected
forest.  After selecting an incoming encounter for each chosen forest
edge, the C components are exactly the isolated collision clusters;
the O forest specifies their blocks.  Consequently these partitions
are recovered from the C/O graph and do not furnish a further partition
factor.  The retained unexpanded independent factors and the continued
labels specify the lower-slab attachment.  This is also why taking a
tree bound for a compatibility coefficient does not remove the
fresh-label denominator.

There are at most \(r\le C_K\) redundant contacts.  Their ordered pair
labels cost at most \((1+n)^{2r}\), and inserting them among the merger
contacts costs at most \((1+n+r)^r\).  Initial forests and the choices
of the bounded exceptional labels have the same polynomial form.
Each ordinary contact has at most two C/O flags, which gives an
absolute exponential base.  A chronological order already determines
the order of contacts on every line; it has no separate permutation
factor.  These observations give
\eqref{eq:sharp-polynomial-count}.  The bounded exceptional counts
alter \(C_K,d_K\), not the base multiplying each ordinary merger.
\end{proof}

\begin{lemma}[Uniform summability]\label{lem:sharp-summability}
For fixed \(K\), the sum of the absolute connected diagram measures
with slab sizes \((n_i)\), after removal of the factor
\(\mu^{1-h}\), is bounded by
\begin{equation}\label{eq:sharp-summable}
       C_K(1+n)^{d_K}(C\tau\sqrt L)^n.
\end{equation}
The same bound, with the appropriate factor \(\mu^{c-h}\), holds
after a bounded number of initial links are erased.  It also holds
after prescribing the absolute spatial position of one observation
in each component, or after inserting a sufficiently small Gaussian
weight in the designated velocities.
\end{lemma}

\begin{proof}
We give the energy and time estimates jointly, to avoid putting a
separate logarithmic velocity loss at each vertex.  Let
\(E_0=\sum_{p=1}^m|v_p^{\mathrm{birth}}|^2\).
Elastic collisions conserve the energy of each colliding pair, O
contacts do not change it, and a line ending at an interface discards
its energy.  Introducing a birth adds the energy already included in
\(E_0\).  Hence the sum of squared velocities of all lines active at
any time is at most \(E_0\).

On the bin \(E_0\in[P/2,P]\), with the bin \(E_0<1\) included by
using \(P=1\), the one-step sum of incoming speed factors over all
choices of a merger pair in slab \(i\) satisfies
\begin{equation}\label{eq:sharp-pair-flux}
 \sum_{p<q}|v_p-v_q|
 \le (n_i-1)\sum_p|v_p|
 \le n_i^{3/2}\sqrt P.
\end{equation}
This estimate is applied to a discrete tree of successive pair
choices at fixed merger times, normals, and birth velocities.  At a
node, all earlier velocities and all intervening specified redundant
contacts are already determined by
\Cref{lem:sharp-merger-coordinates}.  The sum of the weights of its
children is bounded by \eqref{eq:sharp-pair-flux}; an unrealizable
child contributes zero.  Induction over that discrete tree bounds
the sum of the products of fluxes by
\begin{equation}\label{eq:sharp-pair-products}
                  \prod_i(n_i^{3/2}\sqrt P)^{a_i}.
\end{equation}
Integrating the normals adds \(C^{\sum_i a_i}\).  The merger times
are globally ordered within each slab, even if some redundant times
are inserted between them.  Dropping further constraints leaves
time volume at most \(\prod_i\tau^{a_i}/a_i!\).

There remains a common translation for each component.  For any
fixed birth velocities and merger coordinates, write the spatial
arguments of its inputs as \(X+d_p\).  Integrate one chosen input
in \(X\) and use the spatial supremum of every other input.  The
Bol norms in \eqref{eq:sharp-input-bounds} give a factor at most
\(C^m\exp(-bE_0)\).  This bound is independent of the displacements
\(d_p\), which can be large.  To prescribe an observation position
instead, solve \(X=x-d_{\mathrm{obs}}\) and use a spatial supremum
for every input.  The Jacobian is one and the same Gaussian bound
holds uniformly in \(x\).

We spell out the velocity-volume step.  By reserving half of the
Gaussian suppression, for a suitable \(b_1>0\),
\[
 \ind_{E_0\in[P/2,P]}e^{-2b_1E_0}
       \le e^{-b_1P/2}e^{-b_1E_0}.
\]
The integral of the last factor over the \(3m\) birth velocity
coordinates is \((\pi/b_1)^{3m/2}\).  On the lowest bin \(E_0<1\),
omit the dyadic suppression factor and use \(P^{a/2}=1\);
integrating the reserved Gaussian still costs only \(C^m\).
Thus no uncounted power of the energy-bin volume remains.
If \(a=\sum_i a_i\), summing the
dyadic energy bins yields
\begin{equation}\label{eq:sharp-energy-sum}
 \sum_{P=1,2,4,\ldots} e^{-b_1P/2}P^{a/2}
                    \le C^{a+1}(1+a)^{a/2}.
\end{equation}
For example, comparison with the integral of
\(e^{-b_1u/4}(1+u)^{a/2}\) proves this by the gamma integral and
Stirling's inequality.  Its constant depends on \(b_1\), not on
\(K\).

Combining these estimates with \Cref{lem:sharp-factorials}, the
remaining quantity is bounded by a fixed polynomial in \(1+n\)
times
\begin{equation}\label{eq:sharp-before-stirling}
 C^{n+a}\tau^a(1+a)^{a/2}
            \prod_i\frac{n_i^{3a_i/2}}{a_i!n_i!}.
\end{equation}
Set \(\delta_i=n_i-a_i\).  By
\eqref{eq:sharp-line-count}, \(\delta_i\ge0\) and
\(\sum_i\delta_i\le C_K\).  The inequalities
\(n_i!/a_i!\le(1+n_i)^{\delta_i}\) and
\(n_i!\ge(n_i/e)^{n_i}\) imply
\[
 \frac{n_i^{3a_i/2}}{a_i!n_i!}
 \le e^{2n_i}(1+n_i)^{\delta_i}n_i^{-n_i/2},
\]
where zero counts contribute one.  Also
\(\tau^a\le\tau^n\max(1,\tau^{-1})^{C_K}\), and
\((1+a)^{a/2}\le C(1+n)^{C_K}n^{n/2}\), after adjusting the
bounded-size cases.  Finally convexity of \(u\log u\) gives
\begin{equation}\label{eq:sharp-layer-entropy}
               \prod_i n_i^{n_i}\ge(n/L)^n.
\end{equation}
Substitution in \eqref{eq:sharp-before-stirling} proves
\eqref{eq:sharp-summable}.

At a fixed time the sum of squared designated velocities is at
most \(E_0\).  At \(q_0\) record times their total is at most
\(q_0E_0\).  A weight with exponent smaller than
\(b_1/(2\max(1,q_0))\) can therefore be absorbed in the Gaussian
factor reserved above.  This proves the weighted assertions.  The
exponent is chosen so that the fraction of birth Gaussian reserved
for all record times together is fixed; the exponential base
therefore does not grow with the number of observations.  The powers
supplied by initial physical links
and erased links are those already computed in
\eqref{eq:sharp-scaling} and do not affect the summation argument.
\end{proof}

Since \(\tau=T/L\), the base in \eqref{eq:sharp-summable} is
\(CT/\sqrt L\).  Increase \(L\) once so it is less than one.
The number of vectors \((n_i)\) of total size \(n\) is a polynomial
in \(n\), since \(L\) is fixed.  Thus
\eqref{eq:sharp-summable} proves uniform summability and uniform
vanishing of the size tails for every fixed \(K\).  Polynomial
degrees may grow with the factorial order; this never requires
another enlargement of \(L\).

\subsection{The extra-contact null sets}

We next prove the qualitative geometric fact needed for cycles.
Only a fixed finite graph is considered in this subsection.  All
birth velocities, contact times, merger directions, and one free
translation per component are its independent parameters.  They carry
a measure absolutely continuous with respect to the product of
Lebesgue measure and surface measure.  An assertion concerning lines
within the same connected component is invariant under its common
translation and remains valid when one observation position in that
component is fixed as in \Cref{lem:sharp-summability}.  No conditioning
on additional positions or velocities is used in the null-set argument.

\begin{lemma}[No extra contact for a limiting tree]\label{lem:sharp-no-extra}
Consider a finite chronological forest with no physical initial
links, with a prescribed C/O choice on each edge, on its incoming
merger branches.  Outside a null set of its merger parameters the
following hold.
\begin{enumerate}[label=(\roman*)]
\item All specified contacts occur at distinct times away from
interfaces and observation times and are transversal.
\item Before the first redundant contact, the demand for such a
contact is false for all sufficiently small \(\eps\).
\item On a full tree, an additional contact not in the prescribed
list, including a contact preceding the prescribed joining of its
two groups, is absent for all sufficiently small \(\eps\).
\end{enumerate}
In (ii), the independent parameters used are those of the forest
prefix; arbitrary remaining merger parameters may be appended.
In (iii), an outgoing crossing following an O contact is not an
additional contact, because only incoming crossings are counted.
\end{lemma}

\begin{proof}
For clarity first set \(\eps=0\) in the formal trajectories.  At
each prescribed C contact the velocities are transformed by the
linear reflection of the pair velocity in its contact direction;
at an O contact they are unchanged.  Every segment velocity is
therefore polynomial in the birth velocities and in the coordinates
of the sphere directions.  A segment displacement is polynomial
also in the prescribed times.  The parameters lie on the connected
real analytic manifold given by the product of their Euclidean
spaces, the spheres, and an open time simplex.  Consequently a
polynomial expression that is not identically zero has a null zero
set.  It is legitimate to use grazing or ineffective reflections
to test nonidentity: the reflection formulas extend analytically
to those parameters, although the physical flux there may vanish.

There are finitely many pairs of lines and finitely many pairs of
their straight segments.  On a common time interval write their
relative displacement as \(d+tu\).  A coincidence on this interval
requires
\begin{equation}\label{eq:sharp-collinearity}
                             d\times u=0.
\end{equation}
We check the only cases in which this condition could be forced.
If the two lines belong to different forest components, variation
of their free relative translation disproves the identity
\eqref{eq:sharp-collinearity} immediately.  The assertions used for
the first redundant contact and for a connected full tree concern
lines in the same component, as treated next.

Suppose first that the lines are not neighbors in the prescribed
tree.  Make all reflections ineffective.  If the unique joining
path is \(p_0,p_1,\ldots,p_k\), \(k\ge2\), with its edge contact
times \(s_1,\ldots,s_k\), then solving the contact equations gives
the endpoint displacement as a sum of terms
\((v_{p_{j-1}}-v_{p_j})s_j\), plus the endpoint free-flight term.
For \(1\le j<k\), the independent velocity \(v_{p_j}\) occurs
with coefficient \(s_{j+1}-s_j\), up to an immaterial common sign.
These two times are distinct.  Set the endpoint relative velocity
nonzero and vary this intermediate velocity in a direction not
parallel to it.  This disproves the identity
\eqref{eq:sharp-collinearity}.  Branches off the path and birth
interfaces do not alter this test: ineffective reflections leave
their velocities unchanged, and birth times only alter the affine
origins solved for by the contact equations.

Suppose next that the two lines are neighbors and have their
prescribed meeting at time \(a\).  On segments after \(a\), if
neither endpoint has subsequently changed velocity, their flight
is the single separating relative flight from that meeting.  It
cannot give a second incoming crossing, either at zero diameter
or at positive diameter.  This includes an O contact: its relative
flight passes through the sphere and exits once, but does not
enter it a second time without a velocity change.

If an endpoint has a subsequent scattering at time \(b\), take
the first effective such scattering in testing nonidentity and
make all other reflections ineffective.  Its third particle is
distinct from the original pair, since another edge between that
pair would form a cycle.  With all other reflections ineffective,
this third particle's current velocity equals its independent birth
velocity, even if it has a prescribed earlier history.  One can choose
this velocity and the scattering normal so that the relative
velocities of the original pair before and after \(b\) are two
nonparallel vectors \(u_1,u_2\).  Their displacement on the later
segments then contains
\[
                   (b-a)u_1+(t-b)u_2.
\]
Its cross product with \(u_2\) is nonzero, since \(b\ne a\).
For an explicit choice, keep the untouched endpoint at velocity
zero, take the other endpoint's velocities to be \(e_1\) before
and \(e_2\) after the scattering, and use a third-particle velocity
\(-e_1\) with normal \((e_1-e_2)/\sqrt2\).  The equal-mass
reflection gives the stated outgoing endpoint velocity.  This
again disproves \eqref{eq:sharp-collinearity}.  For a segment before
the neighboring pair's prescribed meeting, choose an intervening
scattering time \(b<a\) with the same change from \(e_1\) to
\(e_2\).  Before \(b\), the relative displacement required to
meet at \(a\) is
\[
                 (t-b)e_1-(a-b)e_2,
\]
which is not parallel to the segment relative velocity \(e_1\).
With no intervening reflection the earlier relative flight has
only its prescribed incoming crossing.  This proves the
earlier-segment case without any condition on time-reversed
feasibility.

These arguments use the entire tree and therefore also exclude an
unintended earlier contact between groups whose relative translation
is fixed only by a later prescribed merger.  For a first required
redundant contact, they need only the joining tree in the prefix:
the two lines are already in the same group, and all their relative
parameters are fixed there.  Adding parameters belonging to the
rest of the graph preserves a null exceptional set by Fubini's
theorem.

Zero relative velocities and grazing prescribed contacts have
null parameter sets.  Indeed the incoming relative velocity at a
merger is the difference of velocities from two disjoint previous
groups; it is not an identically zero linear function of their
independent birth velocities.  The normal is a free sphere
parameter.  Specified time coincidences, including equality with
an interface or an observation time, have null time measure.  A
third line at a prescribed vertex has no forced coincidence:
at least one of its paths to the two incident lines is a
nonneighbor path, to which the intermediate-velocity argument
applies.  Endpoints of line lifetimes at interfaces are treated
by the same finite-segment test.  We discard this finite union of
null sets.

It remains to justify that an absent zero-diameter coincidence
excludes positive-diameter contacts, particularly contacts converging
to a prescribed vertex.  For fixed merger times, directions, and
birth velocities the formal tree velocities do not depend on
\(\eps\).  Write each line as
\(x_p(t)=X_p+Y_p(t)\), with \(Y_p\) determined by its velocities.
The edge equation is
\[
 X_p-X_q=-\eps\omega_{pq}-Y_p(s_{pq})+Y_q(s_{pq}).
\]
Solving these equations along the forest shows that, after fixing
its representative translations,
\begin{equation}\label{eq:sharp-uniform-tree-convergence}
        x_p^\eps(t)=x_p^0(t)+\eps d_p,
\end{equation}
where each \(d_p\) is a fixed finite sum of merger normals.
Thus convergence is uniform on every active line interval.

Away from small neighborhoods of the finitely many prescribed
vertices, every pair of closed straight-segment pieces either
has a strictly positive minimal separation or is the already
treated separating flight of a neighboring pair.  For sufficiently
small diameter no new contact can occur there.  Near a prescribed
vertex, a different pair would require an additional zero-offset
pair coincidence, already excluded by the finite-segment argument.
This includes a third-particle coincidence at that vertex and a
disjoint pair coinciding elsewhere at the same time.  The incident
pair has no other velocity
change in a sufficiently small fixed neighborhood, because all
event times are distinct.  It has only the one prescribed incoming
crossing on that neighborhood.  In particular two incoming
crossings cannot collapse onto that vertex.  The same compactness
argument applies to interfaces and observation times.  This proves
(ii) and (iii), including contacts whose times vary with \(\eps\).
\end{proof}

For a graph containing a cycle, keep its first redundant-contact
indicator.  All preceding contacts form a forest, and the relative
trajectory tested by this indicator is determined by that forest
prefix.  No later redundant-contact velocity update enters this test.
By \Cref{lem:sharp-no-extra} it tends to zero for almost
every choice of the independent parameters preceding that contact.
After a finite size truncation, velocities and representative
translations can first be truncated as well.  On those compact
sets the integrands after \Cref{lem:sharp-merger-coordinates} have
an integrable bound, obtained either directly from the finitely
many speed factors or from the proof of
\Cref{lem:sharp-summability}.  One can extend an unrealizable
branch by zero.  Dominated convergence therefore makes the
scaled absolute integral tend to zero.  The velocity and spatial
truncations are then removed by the Gaussian bounds and the uniform
spatial tightness supplied by the Bol-norm convergence in
\eqref{eq:sharp-input-bounds}, and
finally the size truncation is removed using
\eqref{eq:sharp-summable}.  This order of limits needs no rate
uniform in the number of vertices.  Together with
\eqref{eq:sharp-scaling}, it proves the asserted little-o for
cycles and nontrivial initial cumulants.

\subsection{Tree convergence and spatial disintegration}

For a fixed tree all its prescribed C-clusters and O blocks are
below their cutoffs once \(\eps\) is small.  By
\Cref{lem:sharp-no-extra}, almost every set of merger parameters
has only the prescribed incoming contacts for all sufficiently
small \(\eps\).  Thus all history acceptance restrictions then
hold, and counting a selected encounter instead of its existence
changes nothing.  The connected compatibility coefficient on
this configuration is exactly its tree sign \((-1)^{\#O}\).
There is no nontrivial initial link.  The changes of relative
translations therefore leave precisely the fluxes in
\Cref{thm:sharp-cumulants}, after multiplication by
\(\mu^{h-1}\).

Replacing each \(g_i\) by \(f_{i\tau}\) in a fixed-size term is
justified by a telescoping product and
\eqref{eq:sharp-input-bounds}.  The proof of
\Cref{lem:sharp-summability}, with one of the factors replaced by
its Bol-norm difference, gives convergence to zero for that
difference.  For these continuous limiting inputs,
\eqref{eq:sharp-uniform-tree-convergence} and the absence of
contacts at observation times give convergence of every bounded
continuous record test.  The remaining dominated convergence in
spatial translations can be checked first on bounded velocities
and bounded size.  On that set the relevant input displacements
are bounded; the summable cell envelopes in the Bol norm control
the spatial tails uniformly.  Then use
\Cref{lem:sharp-summability} to remove the truncations.  This
defines the limiting signed tree measures.

\begin{lemma}[Spatial traces of tree measures]\label{lem:sharp-spatial-traces}
Let \(\mathcal T\) be any of the absolutely summable tree measures
just constructed, and designate a line and a time at which its
position is required.  It has a version of the disintegration
\begin{equation}\label{eq:sharp-disintegration}
               \mathcal T=\dd x\,\mathcal T_x
\end{equation}
with respect to that position.  For a sufficiently small
\(\alpha>0\), the Gaussian-weighted variation of
\(\mathcal T_x\) is bounded uniformly in \(x\); its integral over
\(x\) is finite.  For every bounded continuous test of the other
coordinates, and also for continuous tests of smaller Gaussian
growth, \(x\mapsto\mathcal T_x(\Phi)\) is continuous.

If two separate tree components are joined at their designated
positions, the product of their conditional measures at a common
position, integrated against incoming flux, is an absolutely
defined measure.  It is the result of imposing the joining
position in the merger coordinates.  Equivalently it is obtained
by a smooth approximate identity for equality of the two
positions.  These facts apply to the marked one-root and
connected two-root tree measures and their record projections.
\end{lemma}

\begin{proof}
For a finite tree choose its free representative translation to
be the position \(x\) of the designated observation.  Every
birth position is then \(x+d_p\), and every recorded position
has the same form, with displacements independent of \(x\).
The derivative in the representative translation is the identity.
This constructs \eqref{eq:sharp-disintegration} as an actual
integral formula, not just as an almost-everywhere equivalence
class.  The fixed-position and integrated versions of
\Cref{lem:sharp-summability} give both asserted variation bounds.

After truncation in size and velocity, the integrand against a
continuous record test is continuous in \(x\), by continuity of
the limiting inputs and the observation updates.  It has the
same integrable majorant uniformly for \(x\) in a compact set.
The velocity tails and the graph-size tails converge uniformly
in \(x\), by the fixed-position estimate.  This proves the
stated continuity.  Gaussian growth in the test is handled by
reserving a slightly stronger Gaussian weight first.

Let the designated variables in two different components have
positions \(x,x_*\).  Their joint integral before joining uses
two independent translations and the product measure.  Setting
\(x_*=x\) gives the claimed conditional product.  The bound
\(B(v-v_*,\omega)\le |v|+|v_*|\), with a slight Gaussian
exponent loss, makes its variation finite: use the uniform
conditional bound for one component and the integrated bound
for the other.  The same estimate dominates a compactly
supported smooth approximate identity in \(x-x_*\), uniformly
in its width.  On bounded velocity sets conditional weak
continuity implies convergence of the product measures against
continuous tests; products of continuous single-component
tests suffice first, and their finite sums approximate a
continuous test on each compact set.  The Gaussian and spatial
variation bounds then remove the compact truncations.  This
proves the approximate-identity characterization and its
agreement with merger coordinates.  Forgetting records commutes
with every displayed integral.
\end{proof}

\subsection{Completeness and splitting of the retained trees}

The preceding size limit was taken with bounded complexity.  It
is necessary to verify that a fixed sufficiently large complexity
cutoff includes every leading tree allowed by the regrouping rule,
even though it may have arbitrarily many ordinary vertices.

\begin{lemma}[Completeness of the tree class]\label{lem:sharp-tree-completeness}
For fixed designated roots, fixed required records, and fixed
\(L\), all global trees with trivial initial cumulants in the
retained expansion after the one-root regrouping of
\Cref{lem:marked-regrouping} have complexity at most a fixed constant
\(K_0\).  This bound is independent of their number of ordinary
contact vertices.  Undoing and reapplying an unmarked one-root
regrouping preserves their weights and their record projections.
\end{lemma}

\begin{proof}
We first prove that an unmarked dependent one-root component
cannot be a tree with trivial initial cumulants.  In its highest
occupied slab, a one-root component has a single connected
block.  If all lower inputs of this block are independent, it
is precisely a term in the definition of \(g_i\), and is
regrouped into that independent factor.  If a connected lower
dependent component attaches to the block through two lines,
the upper connection between these lines and their lower
connection form a global cycle.  For a global tree, every
lower dependent component therefore attaches through at most
one line.  At least one such component must exist for the
upper one to remain dependent.  Following it downward repeats
the same argument on a smaller number of slabs.  At time zero
there is no dependent singleton, and a dependent nonsingleton
requires a nontrivial initial cumulant.  This contradiction
proves the claim by downward induction.  Required records are
the only reason to expand an otherwise independent singleton
history further.

Now cut a global forest at an interface and contract each of
its upper components to a vertex.  Suppose there are \(u\)
such vertices.  A lower connected component with \(k\ge2\)
attachments cannot touch the same upper component twice, since
that creates a cycle.  The bipartite incidence graph of these
lower components and the upper vertices is a forest.  If there
are \(v\) lower vertices of degree at least two and \(d\)
incidence edges, then \(d\le u+v-1\) in each nonempty connected
part, while \(2v\le d\).  In particular
\begin{equation}\label{eq:sharp-crossing-bound}
                               d\le 2u.
\end{equation}
A lower component attached at a single line either is absorbed
into an independent input, by the first paragraph, or continues
a required record.  The latter contribute at most the fixed
number of recorded lines at that interface.  Empty continued
lines are included in this count.

Starting with the fixed number of final roots and iterating
\eqref{eq:sharp-crossing-bound} through \(L\) slabs therefore
bounds the total number of interface roots and recorded
singleton continuations.  There are no cycle or initial-link
contributions to complexity in a leading tree.  This gives
\(K_0\).  The argument applies separately to disconnected
expressions and to projections with fewer required records.

Regrouping an unmarked one-root component is the defining
finite sum for its independent factor, with its original
fresh-label coefficient.  Its components attach at one line
only.  Thus undoing it neither changes a connection between
different remaining components nor introduces a constraint
between them.  The exact finite identity in
\Cref{thm:marked-expansion} shows that forgetting records and
regrouping commute before a limit.  Absolute summability from
\Cref{lem:sharp-summability} gives the same identity for the
limiting trees.  Their weights and projections are therefore
preserved.
\end{proof}

\begin{lemma}[Last-contact products]\label{lem:sharp-tree-splitting}
In the leading tree sums, deleting a last contact of the highest
slab that meets a designated root splits its global tree into
the product of the two component tree measures, with the
incident lines designated before that contact.  Their product
at the collision position is the conditional product in
\Cref{lem:sharp-spatial-traces}.  Conversely, joining such two
disjoint component histories at that contact gives exactly
the corresponding retained tree, with its original factorial
weight.  Record portions not used by a test are projected out.
These operations are valid for bounded continuous record tests
and, with a smaller exponent, for tests of Gaussian velocity
growth.
\end{lemma}

\begin{proof}
Temporarily expand all independent factors in the highest slab by
their one-root definitions for that slab, retaining their lower inputs
in the canonical regrouped form.  Bound the total number of lines in
each resulting limiting tree by \(N\).
For sufficiently small \(\eps\), every particle cap \(\Lambda_i\)
and retained cluster cap \(M_i\) exceeds \(N\).  The convergence
already proved therefore supplies every such leading tree with its
original factorial weight.  We perform the following finite algebra
on this limiting tree class, preserving the inherited bound on the
combined size after a split.  This is not a claim that arbitrary
finite-\(\eps\) capped classes are closed under joining.

Deleting an edge of a global tree makes exactly two components.
Each contains one incident
line of the deleted contact, which is now designated.  If a
component in the highest slab loses its old designated root,
it contains the newly designated incident line or another
original root.  Thus it still satisfies the rooted-block rule.
All lower attachments are unchanged.  They can be put in the
canonical retained form by the singleton regrouping of
\Cref{lem:sharp-tree-completeness}.

Conversely, the two components are disjoint before joining.
Adding their incident edge creates a tree, and its highest
block meets the original root.  The inverse singleton
regroupings restore precisely its retained representation.
This is a bijection at the level of the temporarily expanded
labelled trees.  In particular it is not a factorization of
the actual particle law at an interface.

For the coefficients, choosing an additional incident particle
among \(j\) fresh labels changes \(1/j!\) into
\(1/(j-1)!\).  If the remaining fresh labels are divided into
sets of sizes \(j_1,j_2\), their distribution contributes the
binomial coefficient satisfying
\[
 \binom{j_1+j_2}{j_1}\frac1{(j_1+j_2)!}
                          =\frac1{j_1!j_2!}.
\]
Apply this identity in each slab.  A contact between two
already designated particles uses their single unordered
pair; it has no fresh-label factor.  The sign of the deleted
contact is its C/O sign, and all other signs multiply over
the two components.  This proves the coefficient assertion.

The finite identities retain the combined-size restriction just
specified.  Letting \(N\) increase removes that restriction and
passes the sums through the deletion and joining integrals by
\Cref{lem:sharp-summability,lem:sharp-spatial-traces}.  The
uniform conditional Gaussian bound controls the velocity
flux introduced by the join, with a fixed exponent loss.
Record projection commutes with the finite hierarchy and
with the one-root regrouping, as established above.  Hence
a temporarily designated partner for which no past record
appears in the integrand requires only the projected
component measure.  Approximation locally uniformly with
a stricter Gaussian exponent extends the identity to the
stated test class.
\end{proof}

\begin{proof}[Proof of \Cref{thm:sharp-cumulants}]
We first justify weighted control of the retained high-complexity
tail.  Choose the common birth exponent \(b>0\) in
\eqref{eq:sharp-input-bounds} independently of \(L\), and choose
\(L\) large enough for the coarse package with the fixed smaller
exponent \(3b/4\) as well.  For a fixed observation list let \(J\)
be the number of designated endpoint and record velocity slots.
On every original retained history, conservation of energy gives
\[
 \sum_{\text{these slots}}|v|^2\le J E_0.
\]
Choose the weight exponent \(\alpha>0\) with \(\alpha J\le b/4\).
Its Gaussian weight is dominated by \(e^{bE_0/4}\), which can be
absorbed by replacing every absolute birth input \(|g_i(x,v)|\)
or \(f_0(x,v)\) by that input times \(e^{b|v|^2/4}\).
The resulting inputs have a common \(\Bol,3b/4\) bound.
The positive coarse estimate in \Cref{thm:dhm-package}(a)
requires this envelope, not closeness to the Boltzmann solution.
Apply it to the tilted inputs, with the same initial-link and
fixed-record bounds used in \Cref{thm:marked-expansion}.
Dummy inputs, if used for that bound, may be tilted too;
their additional birth energy only enlarges the majorant.
The same geometric complexity sum makes the weighted retained
tail smaller than any assigned power of \(\eps\), after choosing
\(K\) for the fixed order and accuracy.  The Gaussian fraction
reserved here is fixed, so its atom base does not depend on
\(J,h,K\).  Only the observable weight exponent decreases with
the record list.  This argument concerns retained diagrams,
not a microscopic expansion remainder.

Use \Cref{thm:marked-expansion} with an accuracy
\(\eps^M\) where \(M>2(h-1)\), and with a complexity cutoff
large enough both for that accuracy and for
\Cref{lem:sharp-tree-completeness}.  The expansion error and
the high-complexity tail are negligible after multiplication
by \(\mu^{h-1}\).  They need only their unweighted variation
bound, because the observation tests here are bounded.

For the remaining connected diagrams,
\eqref{eq:sharp-scaling} and
\Cref{lem:sharp-summability} prove the uniform bound at scale
\(\mu^{1-h}\).  Physical initial links gain \(\eps\);
nonphysical initial tails gain any assigned power after
choosing \(D\) as above.  Cycles are negligible by
\Cref{lem:sharp-no-extra}, finite-size dominated convergence,
and the uniform size tails.  The diagrams that remain are
the complete class of leading trees, by
\Cref{lem:sharp-tree-completeness}.  Their fluxes, signs,
input limits, and continuous record tests were evaluated
above.  The same summable majorant proves the asserted
Gaussian variation bounds for their retained sums and
their limits.  In particular it gives the one-root and
two-root limits needed in the next section.  All exponential
bases used in this proof are independent of \(K\) and of
the fixed order \(h\); dependence on those quantities occurs
only in front constants and polynomial degrees.  Thus the
choice of \(\Gamma,L\) remains the single choice made before
the moment order was specified.
\end{proof}
